\begin{example}\label{ex:formM}
Below we give an example of the construction of the matrices $M'$ and $M$.  Let $w = 5237164$, $y = z_{4,6}$, \[
q_a = \det \begin{bmatrix} 0& 0 & {\color{blue}\underline{z_{1,5}}}\\
{\color{blue}\underline{z_{2,2}}} & z_{2,3} & z_{2,5}\\
z_{3,2}& {\color{blue}\underline{z_{3,3}}} & z_{3,5}
 \end{bmatrix} \in Q_{y, I_w}, \mbox{ and } h_b = \det \begin{bmatrix} 
 0 & {\color{blue}\underline{z_{2,2}}} & z_{2,3} & z_{2,4} \\
0 & z_{4,2} & {\color{blue}\underline{z_{4,3}}} & z_{4,4} \\
{ \color{blue} \underline{z_{5,1}}} & z_{5,2} & z_{5,3} & z_{5,4} \\
  z_{6,1} & z_{6,2} & z_{6,3} & {\color{blue}\underline{z_{6,4}}} \\
 \end{bmatrix} \in N_{y,I_w},
 \] in which case $LT(q_a) = z_{1,5}z_{2,2}z_{3,3}$, $LT(h_b) = z_{2,2}z_{4,3}z_{5,1}z_{6,4}$, and \[ 
 LCM(LT(q_a),LT(h_b)) = z_{1,5}z_{2,2}z_{3,3}z_{4,3}z_{5,1}z_{6,4} = LT(\det(M)).
 \]  The variables dividing the leading terms appearing throughout this example are noted in blue and also underlined.  Then \[
 M' = \begin{bmatrix} 
 0 & 0 & 0 & 0 & {\color{blue}\underline{z_{1,5}}}\\
 0 & {\color{blue}\underline{z_{2,2}}} & z_{2,3} & z_{2,4} & z_{2,5}\\
  0 & z_{3,2} & {\color{blue}\underline{z_{3,3}}} & z_{3,4} & z_{3,5}\\
0 & z_{4,2} & {\color{blue}\underline{z_{4,3}}} & z_{4,4} & z_{4,5}\\
  {\color{blue}\underline{z_{5,1}}} & z_{5,2} & z_{5,3} & z_{5,4} & z_{5,5} \\
  z_{6,1} & z_{6,2} & z_{6,3} &  {\color{blue}\underline{z_{6,4}}} & z_{6,5} 
 \end{bmatrix} \mbox{  \hspace{0.5cm}   and  \hspace{0.5cm}   } M  = \begin{bmatrix} 
 0 & 0 & 0 & 0 & 0 & {\color{blue}\underline{z_{1,5}}}\\
 0 & {\color{blue}\underline{z_{2,2}}} & z_{2,3} & z_{2,3} & z_{2,4} & z_{2,5}\\
  0 & z_{3,2} & {\color{blue}\underline{z_{3,3}}} & {\color{blue}\underline{z_{3,3}}} & z_{3,4} & z_{3,5}\\
0 & z_{4,2} & 0 & {\color{blue}\underline{z_{4,3}}}  & z_{4,4} & 0\\
  {\color{blue}\underline{z_{5,1}}} & z_{5,2} & 0 & z_{5,3} & z_{5,4} &0 \\
  z_{6,1} & z_{6,2} & 0  & z_{6,3}& {\color{blue}\underline{z_{6,4}}}  & 0
 \end{bmatrix}.
 \] 

 By expressing $\det(M)$ as a sum of products of $3$-minors in the first $3$ rows of $M$ with $3$-minors in the final $3$ rows, we see $\det(M) \in Q_{y,I_w}$, and, by expressing $\det(M)$ as the sum of products of $4$-minors in columns $1$, $2$, $4$, and $5$ with $2$-minors in columns $3$ and $6$, we see $\det(M) \in N_{y,I_w}$.  Observe that \begin{align*}
 LT(\det(M)) &=  LT\left(\det \begin{bmatrix} z_{2,2} & z_{2,3} \\ z_{3,2} & z_{3,3} \end{bmatrix} \right) \cdot LT\left(\det  \begin{bmatrix} 0 & z_{4,3} & z_{4,4} \\ z_{5,1} & z_{5,3} & z_{5,4} \\ z_{6,1} & z_{6,3} & z_{6,4} \end{bmatrix}\right) \cdot z_{1,5} \cdot\\
 &= (z_{2,2}z_{3,3})(z_{4,3}z_{5,1}z_{6,4})(z_{1,5}).
 \end{align*}
 This expression corresponds to the product $\mu_1 \cdot \mu_2 \cdot \mu_3$ in the proof of Lemma \ref{lem:GrobnerLink}.  One may prefer to use $\widehat{M}$, as in the lemma, obtained from $M$ by subtracting column $3$ from column $4$, which has the effect of setting the copies of $z_{2,3}$ and $z_{3,3}$ in column $4$ to $0$.  
 
One aspect of this example that is typical of the general case is that the intersection of the submatrix of of $M'$ northwest of $z_{6,4}$ (playing the role of $e$) and that northwest of $z_{3,5}$ (from $z_{4,6}$ playing the role of $y$) is a rectangular matrix of indeterminates with $0$'s appearing only in full rows along the top and full columns along the western side of the submatrix.  This arrangement follows from the fact that $5237164$ has no obstruction of Type $1$ and gives rise to the decomposition of $LT(\det(M))$ described above into a product of the leading term of entries along the main diagonal of a matrix of indeterminates together with entries coming only from $q_a$ and entries coming only from $h_b$.
\end{example}
